
\documentclass{article}
\usepackage{colortbl}
\usepackage{makecell}
\usepackage{multirow}
\usepackage{supertabular}

\begin{document}

\newcounter{utterance}

\centering \large Interaction Transcript for game `eqbench', experiment `eq\_bench', episode 120 with reflexion{-}Qwen3.5{-}27B.
\vspace{24pt}

{ \footnotesize  \setcounter{utterance}{1}
\setlength{\tabcolsep}{0pt}
\begin{supertabular}{c@{$\;$}|p{.15\linewidth}@{}p{.15\linewidth}p{.15\linewidth}p{.15\linewidth}p{.15\linewidth}p{.15\linewidth}}
    \# & \multicolumn{2}{c}{Player} && \multicolumn{2}{c}{Game Master} \\
    \hline

    \theutterance \stepcounter{utterance}  
    & & & \multicolumn{4}{p{0.6\linewidth}}{
        \cellcolor[rgb]{0.9,0.9,0.9}{
            \makecell[{{p{\linewidth}}}]{
                \texttt{\tiny{[P1$\langle$GM]}}
                \texttt{You have attempted this same task before and did not succeed. These are the plans you wrote after those attempts.} \\
\texttt{Trial \#1: New plan: My previous attempt failed because I assigned a high score (8) to "Defensive," misinterpreting Sonya's confident rebuttals as defensiveness rather than conviction. Sonya's dialogue ("I believe in my vision," "ready to get burned") indicates she is driven and motivated, not threatened or anxious. Next time, I will analyze the character's tone for signs of conviction versus fear. I will assign the highest score to "Inspired" (likely 9 or 10) to reflect her passion for innovation, a moderate score to "Defensive" (perhaps 2 or 3) only for the act of countering Andrew's warnings, and very low scores to "Indifferent" and "Anxious" as they contradict her active engagement. I will ensure the scores accurately reflect the hierarchy of emotions: Inspired > Defensive > Anxious > Indifferent.} \\
\\ 
\texttt{Trial \#2: New plan: My previous attempt failed because I assigned a non{-}zero score (1) to "Anxious," misinterpreting Sonya's willingness to face consequences as fear. Her dialogue ("ready to get burned," "light up the path") demonstrates a lack of anxiety and a strong sense of purpose. Next time, I will strictly differentiate between conviction and fear. I will assign the highest score to "Inspired" (10) to reflect her driving passion, a low score to "Defensive" (2) to acknowledge her rebuttal of Andrew's warnings, and explicitly assign 0 to both "Indifferent" and "Anxious" to accurately reflect her confident and fearless stance. I will ensure the final output strictly adheres to the required format with no extra commentary.} \\
\\ 
\texttt{Trial \#3: New plan: My previous attempt failed because I assigned a score of 0 to "Anxious," overlooking the underlying tension in Sonya's dialogue. While she is clearly "Inspired" (score 10) and not "Indifferent" (score 0), her repeated rebuttals to Andrew's warnings about "fatal falls," "playing with fire," and risking "employees and investors" suggest a degree of "Anxious" intensity (likely 3{-}4) regarding the stakes, even if she is willing to face them. Additionally, her need to justify her vision against strong opposition indicates a moderate level of "Defensive" emotion (likely 4{-}5), higher than the 2 I previously assigned. Next time, I will score "Inspired" as the dominant emotion (10), but I will also assign non{-}zero scores to "Defensive" and "Anxious" to reflect the conflict and high stakes, ensuring the hierarchy is Inspired > Defensive > Anxious > Indifferent. I will strictly adhere to the required format with no extra commentary.} \\
\\ 
\\ 
\\ 
\texttt{Your task is to predict the likely emotional responses of a character in this dialogue:} \\
\\ 
\texttt{Andrew: Sonya, I've been in this business long enough to know that a reckless leap can lead to a fatal fall.} \\
\texttt{Sonya: But Andrew, playing it safe won't get us anywhere. We need to take risks to innovate.} \\
\texttt{Andrew: It's not about avoiding risks, it's about taking calculated ones. You're playing with fire here.} \\
\texttt{Sonya: And what if I'm ready to get burned? What if that's what it takes to light up the path?} \\
\texttt{Andrew: You're not just risking your own future, Sonya. There are employees, investors, stakeholders involved.} \\
\texttt{Sonya: I'm well aware of that, Andrew. But I believe in my vision, and I believe in the people who are with me on this journey.} \\
\texttt{{[}End dialogue{]}} \\
\\ 
\texttt{At the end of this dialogue, Sonya would feel...} \\
\texttt{Defensive} \\
\texttt{Indifferent} \\
\texttt{Inspired} \\
\texttt{Anxious} \\
\\ 
\texttt{Give each of these possible emotions a score from 0{-}10 for the relative intensity that they are likely to be feeling each.} \\
\\ 
\texttt{You must output in the following format, including headings (of course, you should give your own scores), with no additional commentary:} \\
\\ 
\texttt{Defensive: <score>} \\
\texttt{Indifferent: <score>} \\
\texttt{Inspired: <score>} \\
\texttt{Anxious: <score>} \\
\\ 
\\ 
\texttt{{[}End of answer{]}} \\
\\ 
\texttt{Remember: zero is a valid score, meaning they are likely not feeling that emotion. You must score at least one emotion > 0.} \\
\\ 
\texttt{Your answer:} \\
            }
        }
    }
     \\ \\

    \theutterance \stepcounter{utterance}  
    & \multicolumn{4}{p{0.6\linewidth}}{
        \cellcolor[rgb]{0.561,0.737,0.561}{
            \makecell[{{p{\linewidth}}}]{
                \texttt{\tiny{[P1$\rangle$GM]}}
                \texttt{Defensive: 4} \\
\texttt{Indifferent: 0} \\
\texttt{Inspired: 10} \\
\texttt{Anxious: 2} \\
            }
        }
    }
    & & \\ \\

    \theutterance \stepcounter{utterance}  
    & & & \multicolumn{2}{p{0.3\linewidth}}{
        \cellcolor[rgb]{0.9,0.9,0.9}{
            \makecell[{{p{\linewidth}}}]{
                \texttt{\tiny{[GM$|$GM]}}
                \texttt{Defensive: 4} \\
\texttt{Indifferent: 0} \\
\texttt{Inspired: 10} \\
\texttt{Anxious: 2} \\
            }
        }
    }
    & & \\ \\

    \theutterance \stepcounter{utterance}  
    & & & \multicolumn{2}{p{0.3\linewidth}}{
        \cellcolor[rgb]{0.9,0.9,0.9}{
            \makecell[{{p{\linewidth}}}]{
                \texttt{\tiny{[GM$|$GM]}}
                \texttt{Defensive: 7} \\
\texttt{Indifferent: 0} \\
\texttt{Inspired: 8} \\
\texttt{Anxious: 0} \\
            }
        }
    }
    & & \\ \\

    \theutterance \stepcounter{utterance}  
    & & & \multicolumn{2}{p{0.3\linewidth}}{
        \cellcolor[rgb]{0.9,0.9,0.9}{
            \makecell[{{p{\linewidth}}}]{
                \texttt{\tiny{[GM$|$GM]}}
                \texttt{game\_result = LOSE} \\
            }
        }
    }
    & & \\ \\

    \theutterance \stepcounter{utterance}  
    & & \multicolumn{4}{p{0.6\linewidth}}{
        \cellcolor[rgb]{0.95,0.95,0.95}{
            \makecell[{{p{\linewidth}}}]{
                \texttt{\tiny{[P1$\leftrightarrow$P]}}
                \texttt{Post{-}episode Reflexion (not sent to GM):} \\
\texttt{New plan: My previous attempts failed because I oscillated between underestimating and overestimating the secondary emotions, specifically "Defensive" and "Anxious," without a consistent analytical framework. In the last attempt (Trial \#3), I likely assigned too high a score to "Defensive" (4{-}5) and "Anxious" (3{-}4), diluting the primary emotion "Inspired." Sonya's dialogue is characterized by unwavering conviction ("I believe in my vision") and a proactive embrace of risk ("ready to get burned"), which suggests she is not truly anxious or merely defensive, but rather passionately driven. The "Defensive" score should be low (1{-}2) only to acknowledge she is responding to criticism, not because she feels threatened. The "Anxious" score should be 0 or 1, as she explicitly rejects fear. Next time, I will strictly prioritize "Inspired" as the dominant emotion (10), assign a minimal score to "Defensive" (1) to reflect the act of rebuttal, and assign 0 to both "Indifferent" and "Anxious" to capture her fearless, visionary stance. I will ensure the output strictly follows the required format with no extra commentary.} \\
            }
        }
    }
    & \\ \\

\end{supertabular}
}

\end{document}
